% SkyScore - Báo cáo dự án
% Biên dịch: pdfLaTeX (hoặc XeLaTeX) trên Overleaf / TeX Live / MiKTeX.
\documentclass[12pt,a4paper]{article}

\usepackage[utf8]{inputenc}
\usepackage[T5]{fontenc}
\usepackage[vietnamese]{babel}
\usepackage{lmodern}
\usepackage[margin=2.5cm]{geometry}
\usepackage{amsmath,amssymb,bm}
\usepackage{booktabs,tabularx,array}
\usepackage{graphicx}
\usepackage{xcolor}
\usepackage{listings}
\usepackage{tikz}
\usetikzlibrary{arrows.meta,positioning,shapes.geometric}
\usepackage{enumitem}
\usepackage[hidelinks]{hyperref}

\definecolor{codebg}{RGB}{246,246,246}
\lstset{
  basicstyle=\ttfamily\small,
  backgroundcolor=\color{codebg},
  frame=single, rulecolor=\color{gray!40},
  breaklines=true, columns=fullflexible,
  keywordstyle=\color{blue!60!black}\bfseries,
  commentstyle=\color{gray}\itshape,
  showstringspaces=false
}

\newcommand{\ind}{\mathbb{1}}
\newcommand{\arcsec}{\ensuremath{''}}

\title{\textbf{SkyScore}\\[4pt]
\large Hệ thống lập kế hoạch thời điểm quan sát thiên văn\\
đa mục tiêu kết hợp học máy dự báo seeing}
\author{Nhóm thực hiện: \underline{\hspace{6cm}}}
\date{Tháng 11 năm 2026}

\begin{document}
\maketitle

\begin{abstract}
\noindent
SkyScore trả lời câu hỏi thực tiễn của các câu lạc bộ thiên văn: \emph{``Tối nay, vào lúc nào
nên quan sát một mục tiêu cho trước?''}. Hệ thống kết hợp (i) tính toán vị trí thiên thể bằng
Astropy, (ii) dự báo thời tiết theo giờ từ Open-Meteo, (iii) mô hình Random Forest ước lượng
\emph{seeing} khí quyển và (iv) một hàm điểm đa mục tiêu minh bạch có ràng buộc cứng. Ngoài
khung giờ tốt nhất, hệ thống còn cung cấp tập nghiệm Pareto và phân tích độ nhạy theo trọng số
bằng phân phối Dirichlet. Kết quả được trình bày trên giao diện web Streamlit. Mô hình seeing
hiện được huấn luyện trên nhãn giả lập (MAE $0{,}17\arcsec$, $R^2 = 0{,}96$ theo kiểm định
chéo chuỗi thời gian); chúng tôi thảo luận rõ giới hạn này và lộ trình thay bằng dữ liệu đo thật.
\end{abstract}

\tableofcontents
\newpage

%=====================================================================
\section{Đặt vấn đề}

Chất lượng một buổi quan sát thiên văn phụ thuộc đồng thời vào nhiều yếu tố:
\begin{itemize}[nosep]
  \item \textbf{Hình học thiên văn}: mục tiêu phải đủ cao trên đường chân trời; trời phải đủ tối.
  \item \textbf{Mặt Trăng}: trăng sáng, đã mọc và nằm gần mục tiêu làm bầu trời sáng lên.
  \item \textbf{Thời tiết}: mây che là yếu tố loại trừ gần như tuyệt đối.
  \item \textbf{Seeing}: nhiễu loạn khí quyển làm ảnh sao bị nhoè, đo bằng giây cung (\arcsec).
\end{itemize}
Các yếu tố này có đơn vị khác nhau (độ, phần trăm, giây cung) và thường \emph{mâu thuẫn}
nhau: giờ mục tiêu lên cao nhất có thể trùng lúc trăng sáng hoặc nhiều mây. Bài toán vì vậy
mang bản chất \textbf{tối ưu đa mục tiêu}.

\paragraph{Mục tiêu dự án.}
\begin{enumerate}[nosep]
  \item Xây dựng pipeline tự động từ toạ độ, mục tiêu, ngày $\rightarrow$ bảng đánh giá từng giờ.
  \item Thiết kế hàm điểm minh bạch, có thể giải thích và điều chỉnh theo nhu cầu người dùng.
  \item Ước lượng seeing bằng học máy với quy trình kiểm định không rò rỉ thông tin tương lai.
  \item Đánh giá độ tin cậy của khuyến nghị (Pareto, độ nhạy trọng số).
  \item Cung cấp giao diện web dễ dùng cho thành viên CLB.
\end{enumerate}

%=====================================================================
\section{Kiến trúc hệ thống}

\begin{figure}[h]
\centering
\begin{tikzpicture}[
  node distance=8mm and 12mm,
  box/.style={draw, rounded corners, align=center, minimum height=9mm, inner sep=4pt, fill=blue!5},
  arr/.style={-{Stealth}, thick}
]
\node[box] (astro) {Astropy\\ephemeris};
\node[box, below=of astro] (wx) {Open-Meteo\\API};
\node[box, right=of wx] (feat) {Đặc trưng\\thời tiết};
\node[box, right=of feat] (rf) {Random Forest\\(seeing)};
\node[box, right=22mm of astro] (tbl) {Bảng\\theo giờ};
\node[box, right=of tbl] (norm) {Chuẩn hoá\\$g_i\in[0,1]$};
\node[box, right=of norm] (S) {Điểm $S(t)$};
\node[box, above right=4mm and 10mm of S] (win) {Cửa sổ\\tốt nhất};
\node[box, below right=4mm and 10mm of S] (par) {Pareto,\\độ nhạy};
\draw[arr] (astro) -- (tbl);
\draw[arr] (wx) -- (tbl);
\draw[arr] (wx) -- (feat);
\draw[arr] (feat) -- (rf);
\draw[arr] (rf) -- (norm);
\draw[arr] (tbl) -- (norm);
\draw[arr] (norm) -- (S);
\draw[arr] (S) -- (win);
\draw[arr] (S) -- (par);
\end{tikzpicture}
\caption{Luồng xử lý dữ liệu của SkyScore.}
\end{figure}

Mã nguồn được tổ chức thành các mô-đun độc lập (Bảng~\ref{tab:modules}), mỗi mô-đun có
kiểm thử đơn vị bằng \texttt{pytest}.

\begin{table}[h]
\centering
\begin{tabularx}{\textwidth}{>{\ttfamily}l X}
\toprule
\normalfont\textbf{Tệp} & \textbf{Chức năng} \\
\midrule
astro.py        & Độ cao mục tiêu, Mặt Trời, Mặt Trăng; pha trăng và góc tách trăng--mục tiêu \\
weather.py      & Gọi API Open-Meteo; sinh thời tiết giả lập khi offline \\
seeing\_model.py & Trích đặc trưng, nhãn giả lập, Random Forest, kiểm định chéo chuỗi thời gian \\
scoring.py      & Chuẩn hoá, hàm điểm, cửa sổ tốt nhất, Pareto, độ nhạy \\
pipeline.py     & Ghép nối các bước trên thành một lời gọi \texttt{build\_table} \\
app.py          & Giao diện Streamlit \\
\bottomrule
\end{tabularx}
\caption{Các mô-đun chính.}
\label{tab:modules}
\end{table}

%=====================================================================
\section{Dữ liệu}

\subsection{Hình học thiên văn}
Với vị trí quan sát $(\varphi,\lambda)$ và lưới thời gian 24 giờ (UTC), Astropy tính:
độ cao mục tiêu $h$, độ cao Mặt Trời $h_\odot$, độ cao Mặt Trăng $h_{\text{moon}}$, góc tách
$\theta$ giữa Mặt Trăng và mục tiêu, và phần đĩa trăng được chiếu sáng
\begin{equation}
k = \frac{1+\cos i}{2},
\end{equation}
trong đó $i$ là góc pha. Mục tiêu gồm 6 thiên thể sâu (M31, M42, M45, M13, M8, M104 --- toạ độ
J2000) và 4 hành tinh (Sao Mộc, Sao Thổ, Sao Hoả, Sao Kim --- lịch thiên văn tích hợp của
Astropy, không cần mạng).

\subsection{Thời tiết}
Biến theo giờ từ Open-Meteo (miễn phí, không cần khoá API): độ che phủ mây, độ ẩm tương đối,
nhiệt độ, điểm sương, áp suất mực biển, gió 10\,m và gió ở mực 250\,hPa (dòng xiết).
Khi không có mạng, hệ thống dùng thời tiết giả lập gồm chu kỳ ngày đêm cộng nhiễu AR(1) và
\textbf{ghi rõ nguồn trên giao diện}.

%=====================================================================
\section{Mô hình toán học}

\subsection{Chuẩn hoá tiêu chí}
Mỗi tiêu chí được ánh xạ về \emph{độ tốt} $g_i\in[0,1]$ (1 là tốt nhất):
\begin{align}
g_{\text{alt}}  &= \max(0,\sin h), \label{eq:alt}\\
g_{\text{moon}} &= 1 - k\,\ind[h_{\text{moon}}>0]\left(0{,}4 + 0{,}6\,e^{-\theta/40^\circ}\right), \label{eq:moon}\\
g_{\text{cloud}} &= 1 - c/100, \label{eq:cloud}\\
g_{\text{see}}  &= 1 - \operatorname{clip}\!\left(\frac{\sigma-0{,}8}{3{,}5-0{,}8},\,0,\,1\right). \label{eq:see}
\end{align}
Trong đó $c$ là độ che phủ mây (\%) và $\sigma$ là seeing (\arcsec).
Biểu thức \eqref{eq:alt} xấp xỉ $1/\text{airmass}$ (với $X\approx 1/\sin h$).
Biểu thức \eqref{eq:moon} phạt nặng nhất khi trăng tròn, đã mọc và nằm sát mục tiêu; ngay cả khi
ở xa, trăng sáng vẫn làm sáng nền trời (hệ số nền $0{,}4$).

\subsection{Hàm điểm và ràng buộc}
\begin{equation}
S(t) = \ind[\text{obs}(t)]\sum_{i\in\{\text{alt,moon,cloud,see}\}} w_i\,g_i(t),
\qquad \sum_i w_i = 1,\ w_i\ge 0,
\end{equation}
với ràng buộc cứng
\begin{equation}
\text{obs}(t) \iff h_\odot(t) < -18^\circ \ \wedge\ h(t) > 15^\circ .
\end{equation}
Điều kiện thứ nhất là đêm thiên văn; điều kiện thứ hai loại bỏ vùng airmass lớn ($X>3{,}9$).
Trọng số mặc định: $\bm{w} = (0{,}35;\ 0{,}20;\ 0{,}30;\ 0{,}15)$, người dùng có thể thay đổi
và hệ thống tự chuẩn hoá về tổng bằng 1.

\subsection{Cửa sổ quan sát tốt nhất}
Với độ dài cửa sổ $L$ giờ:
\begin{equation}
t^\ast = \arg\max_{t}\ \frac{1}{L}\sum_{j=0}^{L-1} S(t+j).
\end{equation}

\subsection{Tập Pareto}
Giờ $a$ \emph{trội hơn} giờ $b$ nếu $g_i(a)\ge g_i(b)\ \forall i$ và bất đẳng thức chặt ở ít
nhất một $i$. Tập Pareto là các giờ quan sát được không bị giờ nào trội hơn. Tập này cho thấy
sự đánh đổi giữa các tiêu chí \emph{mà không cần chọn trọng số}.

\subsection{Phân tích độ nhạy trọng số}
Trọng số là lựa chọn chủ quan. Để kiểm tra khuyến nghị có phụ thuộc vào lựa chọn đó không, ta
lấy mẫu $n=500$ bộ trọng số $\bm{w}^{(m)}\sim\operatorname{Dirichlet}(1,1,1,1)$ (phân phối đều
trên đơn hình) và tính tần suất thắng
\begin{equation}
f(t) = \frac{1}{n}\sum_{m=1}^{n}\ind\!\left[t = \arg\max_{\tau} S_{\bm{w}^{(m)}}(\tau)\right].
\end{equation}
$\max_t f(t)$ cao nghĩa là khuyến nghị \textbf{bền vững}.

%=====================================================================
\section{Mô hình học máy dự báo seeing}

\subsection{Đặc trưng}
Seeing chủ yếu do nhiễu loạn tầng cao (dòng xiết), gió mặt đất và biến đổi nhiệt. Bảy đặc trưng:
\begin{itemize}[nosep]
  \item tốc độ gió 250\,hPa, tốc độ gió 10\,m;
  \item $|\Delta T_{3h}|$, $|\Delta p_{3h}|$ (biến thiên nhiệt độ, áp suất trong 3 giờ);
  \item chênh lệch nhiệt độ -- điểm sương, độ ẩm tương đối, độ che phủ mây.
\end{itemize}

\subsection{Nhãn và mô hình}
Do không có kho dữ liệu seeing đo thật miễn phí cho địa điểm cụ thể, nhãn được sinh bằng công
thức gần đúng mang tính vật lý:
\begin{equation}
\sigma = \Big[0{,}7 + 0{,}012\,(v_{250}-40)^+ + 0{,}03\,v_{10} + 0{,}25\,|\Delta T_{3h}|
 + 0{,}15\,|\Delta p_{3h}| + 0{,}004\,(RH-60)^+\Big]\cdot\varepsilon,
\end{equation}
với $\varepsilon\sim\operatorname{LogNormal}(0;\,0{,}08)$, sau đó cắt về $[0{,}5;\,5]\arcsec$.
Tập huấn luyện gồm $24\times240 = 5760$ giờ thời tiết giả lập.

Mô hình: \texttt{RandomForestRegressor} (200 cây, \texttt{min\_samples\_leaf}=3).
Kiểm định: \texttt{TimeSeriesSplit} 5 lần gập --- luôn huấn luyện trên quá khứ, kiểm tra trên
tương lai, tránh rò rỉ thông tin.

\subsection{Kết quả}
\begin{table}[h]
\centering
\begin{tabular}{lc}
\toprule
\textbf{Chỉ số (CV chuỗi thời gian, 5 lần)} & \textbf{Giá trị} \\
\midrule
MAE  & $0{,}17\arcsec$ \\
$R^2$ & $0{,}96$ \\
\bottomrule
\end{tabular}
\quad
\begin{tabular}{lc}
\toprule
\textbf{Đặc trưng} & \textbf{Độ quan trọng} \\
\midrule
Gió 250\,hPa & 0,880 \\
$|\Delta T_{3h}|$ & 0,052 \\
Gió 10\,m & 0,043 \\
$|\Delta p_{3h}|$ & 0,017 \\
Còn lại & $<0{,}01$ \\
\bottomrule
\end{tabular}
\caption{Hiệu năng và độ quan trọng đặc trưng của mô hình seeing.}
\end{table}

\paragraph{Diễn giải trung thực.}
Các con số trên chỉ cho thấy pipeline \emph{khôi phục được công thức giả lập}; chúng
\textbf{không} phải bằng chứng về khả năng dự báo seeing thực tế. Độ quan trọng của gió
250\,hPa ($\approx 0{,}88$) phản ánh cách xây dựng nhãn giả lập.

%=====================================================================
\section{Kết quả minh hoạ}

Ví dụ: TP.\,Hồ Chí Minh ($10{,}78^\circ$N, $106{,}70^\circ$E), mục tiêu M42 (Tinh vân Lạp Hộ),
ngày 15/01/2026 (UTC), trọng số mặc định, chế độ thời tiết giả lập, cửa sổ 3 giờ.

\begin{table}[h]
\centering
\begin{tabular}{lc}
\toprule
\textbf{Kết quả} & \textbf{Giá trị} \\
\midrule
Số giờ quan sát được & 7 \\
Cửa sổ tốt nhất (UTC) & 13:00 -- 16:00 \\
Cửa sổ tốt nhất (giờ Việt Nam) & 20:00 -- 23:00 \\
Điểm trung bình cửa sổ & 0,76 \\
Độ bền vững giờ tốt nhất & 46\,\% \\
Số giờ thuộc tập Pareto & 7 / 7 \\
\bottomrule
\end{tabular}
\caption{Kết quả minh hoạ (dữ liệu thời tiết giả lập).}
\end{table}

Nhận xét: cả 7 giờ đều thuộc tập Pareto, nghĩa là không giờ nào tốt hơn tuyệt đối ở mọi tiêu
chí --- đúng là có sự đánh đổi (thường giữa độ cao mục tiêu và mây). Độ bền vững 46\,\% cho thấy
giờ tốt nhất phụ thuộc đáng kể vào trọng số, nên người dùng nên xem cả cửa sổ nhiều giờ thay vì
một giờ duy nhất.

% Gợi ý: chụp màn hình giao diện và chèn tại đây
% \begin{figure}[h]\centering
% \includegraphics[width=\textwidth]{screenshot_planner.png}
% \caption{Giao diện tab Planner.}\end{figure}

%=====================================================================
\section{Giao diện người dùng}

Ứng dụng Streamlit gồm thanh bên (toạ độ, mục tiêu, ngày, trọng số, độ dài cửa sổ) và ba tab:
\begin{itemize}[nosep]
  \item \textbf{Planner}: cửa sổ tốt nhất, điểm trung bình, độ bền vững; biểu đồ điểm, bốn tiêu
        chí và seeing theo giờ; bảng Pareto và bảng dữ liệu thô.
  \item \textbf{Giới thiệu}: mô tả phương pháp và lưu ý về dữ liệu.
  \item \textbf{Hướng dẫn}: cách dùng, cách đọc kết quả, câu hỏi thường gặp.
\end{itemize}

Chạy cục bộ:
\begin{lstlisting}[language=bash]
pip install -r requirements.txt
pytest                  # kiem thu don vi
streamlit run app.py    # mo giao dien web
\end{lstlisting}

%=====================================================================
\section{Kiểm thử}
Bộ kiểm thử \texttt{pytest} bao gồm: tính đúng của chuẩn hoá min--max, ràng buộc quan sát
được, điểm nằm trong $[0,1]$, cửa sổ tốt nhất, tập Pareto, độ nhạy trọng số, và chạy toàn bộ
pipeline ở chế độ offline.

%=====================================================================
\section{Hạn chế và hướng phát triển}
\begin{enumerate}[nosep]
  \item \textbf{Nhãn seeing thật}: thay nhãn giả lập bằng nhật ký quan sát của CLB (ước lượng
        seeing qua tách sao đôi hoặc FWHM sao trên ảnh), ghép với thời tiết lưu trữ cùng giờ,
        và so sánh với mô hình cơ sở (persistence / climatology).
  \item \textbf{Ô nhiễm ánh sáng}: thêm tiêu chí thứ năm theo thang Bortle hoặc dữ liệu vệ tinh VIIRS.
  \item \textbf{Hiệu chỉnh trọng số}: học trọng số từ đánh giá chủ quan của thành viên sau mỗi buổi quan sát.
  \item \textbf{Múi giờ địa phương}: hiển thị trực tiếp giờ địa phương thay vì UTC.
\end{enumerate}

%=====================================================================
\section{Kết luận}
SkyScore là công cụ lập kế hoạch quan sát minh bạch, kết hợp tính toán thiên văn, dự báo thời
tiết và học máy. Thay vì chỉ đưa ra một con số, hệ thống còn cho biết \emph{vì sao} (bốn tiêu
chí), \emph{có đánh đổi gì} (Pareto) và \emph{chắc chắn đến đâu} (độ nhạy trọng số). Bước quan
trọng nhất tiếp theo là thu thập dữ liệu seeing thật để mô hình học máy có giá trị khoa học.

%=====================================================================
\begin{thebibliography}{9}
\bibitem{astropy} Astropy Collaboration (2022). \emph{The Astropy Project: Sustaining and
  Growing a Community-oriented Open-source Project and the Latest Major Release (v5.0) of the
  Core Package}. The Astrophysical Journal, 935(2), 167.
\bibitem{openmeteo} Open-Meteo. \emph{Free Weather API}. \url{https://open-meteo.com/}
\bibitem{sklearn} Pedregosa, F.\ et al.\ (2011). \emph{Scikit-learn: Machine Learning in
  Python}. Journal of Machine Learning Research, 12, 2825--2830.
\bibitem{breiman} Breiman, L.\ (2001). \emph{Random Forests}. Machine Learning, 45(1), 5--32.
\bibitem{coulman} Coulman, C.\,E.\ (1985). \emph{Fundamental and Applied Aspects of
  Astronomical Seeing}. Annual Review of Astronomy and Astrophysics, 23, 19--57.
\bibitem{streamlit} Streamlit. \url{https://streamlit.io/}
\end{thebibliography}

\end{document}
